% Ad Familiares 8.8 (ad-familiares-08-08)
% Source: https://raw.githubusercontent.com/PerseusDL/canonical-latinLit/master/data/phi0474/phi056/phi0474.phi056.perseus-lat1.xml
% Fetched by scripts/fetch_latin.py

% Dateline (Perseus): Scr. Romae in. m. Oct. a. 703 (51) .

\ciceroLetterOpener{CAELIVS CICERONI S.}

\ciceroSection{1} etsi de re p. quae tibi scribam habeo, tamen nihil quod magis gavisurum te putem habeo quam hoc: scito C. Sempronium Rufum, mel ac delicias tuas, calumniam maximo plausu tulisse. qua quaeris in causa. M. Tuccium, accusatorem silum, post ludos Romanos reum lege Plotia de vi fecit hoc consilio, quod videbat, si extraordinarius reus nemo accessisset, sibi hoc anno causam esse dicendam. Dubium porro illi non erat quid futurum esset. nemini hoc deferre munusculum maluit quam suo accusatori; itaque sine ullo subscriptore descendit et Tuccium reum fecit. at ego, simul atque audivi, invocatus ad subsellia rei occurro; surgo neque verbum de re facio; totum Sempronium usque eo perago ut Vestorium quoque interponam et illam fabulam narrem, quem ad modum tibi pro beneficio dederit, †si quod iniuriis suis esset, ut Vestorius teneret.

\ciceroSection{2} haec quoque magna nunc contentio forum tenet: M. Servilius postquam, ut coeperat, omnibus in rebus turbarat nec quod non venderet quicquam reliquerat maximaeque nobis traditus erat invidiae, neque Laterensis praetor postulante Pausania nobis patronis QVO EA PECVNIA PERVENISSET recipere voluit, Q. Pilius, necessarius Attici nostri, de repetundis eum postulavit. Magna ilico fama surrexit et de damnatione ferventer loqui est coeptum. quo vento proicitur Appius minor, ut †inpicet depecuniam ex bonis patris pervenisse ad Servilium praevaricationisque causa diceret depositum HS LXXXI. admiraris amentiam immo, si actionem stultissimasque de se, nefarias de patre confessiones audisses.

\ciceroSection{3} mittit in consilium eosdem illos qui litis aestimarant iudices. Cum aequo numero sententiae fuissent, Laterensis leges ignorans pronuntiavit quid singuli ordines iudicassent, et ad extremum, ut solent, 'NON REDIGAM. Postquam discessit et pro absoluto Servilius haberi coeptus legisque unum et centesimum caput legit, in quo ita erat: QVOD EORVM IVDICVM MAIOR PARS IVDICARIT, ID IVS RATVMQVE ESTO, in tabulas absolutum non rettulit, ordinum iudicia perscripsit; postulante rusus Appio cum L. Lollio transegit et se relaturum dixit. sic nunc neque absolutus neque damnatus Servilius de repetundis saucius Pilio tradetur. nam de divinatione Appius, cum calumniam iurasset, contendere ausus non est Pilioque cessit et ipse de pecuniis repetundis a Serviliis est postulatus et praeterea de vi reus a quodam suo emissario, S. Tettio, factus. recte hoc par habet.

\ciceroSection{4} quod ad rem publicam pertinet, omnino multis diebus exspectatione Galliarum actum nihil est; aliquando tamen saepe re dilata et graviter acta et plane perspecta Cn. Pompei voluntate in eam partem ut eum decedere post K. Martias placeret, senatus consultum, quod tibi misi, factum est auctoritatesque perscriptae.

\ciceroSection{5} senatus consultum, auctoritates . Pr. K. Octobris in aede Apollinis scrib. adfuerunt L. Domitius Cn. f. Fab. Ahenobarbus, Q. Caecilius Q. f. Fab. Metellus Pius Scipio, L. Villius L. f. pom . annalis , C. Septimius T. f. qui ., C. Lucilius C. f. Pup. Hirrus, C. Scribonius C. f. Pop. Curio, L. Ateius L. f. an . Capito, M. Eppius M. f. ter . quod M. Marcellus cos. v.f. de provinciis consularibus, d. e. r. l. c., uti L. Paulus, C. Marcellus coss., cum magistratum inissent, e x K. Mart. , quae in suo magistratu futurae essent, de consularibus provinciis ad senatum referrent, neve quid prius e x K. Mart. ad senatum referrent neve quid coniunctim de ea re referretur a consiliis , utique eius rei causa per dies comitialis senatum haberent senatusque cons. facerent et, cum de ea re ad senatum referretur, a consiliis, qui eorum in CCC iudicibus essent, s.f. s. adducere liceret; si quid de ea re ad populum pl. ve lato opus esset, uti Ser. Sulpicius, AL Marcellus coss., praetores tr. q. pl., quibus eorum videretur, ad populum pl. ve ferrent;

\ciceroSection{6} quod si ii non tulissent, uti quicumque deinceps essent ad populum pl. ve ferrent C. Pr. K. Octobris in aede Apollinis scrib. adfuerunt L. Domitius Cn. f. Fab. Ahenobarbus, Q. Caecilius Q. f. Fab. Metellus Pius Scipio, L. Villius L. f. pom . annalis , C. Septimius T. f. qui ., C. Lucilius C. f. Pup. Hirrus, C. Scribonius C. f. Pop. Curio, L. Ateius L .f. an . Capito, M. Eppius M.f. ter . quod M. Marcellus cos. v. f. de provinciis, d. e. r. l. c., senatum existimare neminem eorum, qui potestatem habent intercedendi, impediendi, moram adferre oportere quo minus de r. p. p. R. q. p. ad senatum referri senatique consultum fieri possit; qui impedierit, prohibuerit, eum senatum existimare contra rem publicam fecisse. si quis huic s. c. intercesserit, senatui placere auctoritatem perscribi et de ea re ad senatum p. q. t. referri. huic s. c. intercessit C. Caelius, L. Vinicius, P. Cornelius, C. Vibius Pansa, tr. pl.

\ciceroSection{7} item senatui placere de militibus, qui in exercitu C. Caesaris sunt, qui eorum stipendia emerita aut causas, quibus de causis missi fieri debeant, habeant, ad hunc ordinem referri, ut eorum ratio habeatur causaeque cognoscantur. si quis huic s. c. intercessisset, senatui placere auctoritatem perscribi et de ea re p. q. t. adhunc ordinem referri. huic s. c. intercessit C. Caelius, C. Pansa, tr. pl.

\ciceroSection{8} itemque senatui placere in Ciliciam provinciam, in VIII reliquas provincias, quas praetorii pro praetore obtinerent, eos, qui praetores fuerunt neque in provincia cum imperio fuerunt, quos eorum ex s. c. cum imperio in provincias pro praetore mitti oporteret, eos sortito in provincias mitti placere; si ex eo numero, quos ex s. c. in provincias ire oporteret, ad numerum non essent qui in eas provincias proficiscerentur, tum, uti quodque conlegium primum praetorum fuisset neque in provincias profecti essent, ita sorte in provincias proficiscerentur; si ii ad numerum non essent, tum deinceps proximi cuiusque conlegii qui praetores fuissent neque in provincias profecti essent, in sortem coicerentur, quoad is numerus effectus esset, quem ad numerum in provincias mitti oporteret si quis huic s. c. intercessisset, auctoritas perscriberetur. huic s. c. intercessit C. Caelius, C. Pansa, tr. pl.

\ciceroSection{9} illa praeterea Cn. Pompei sunt animadversa, quae maxime confidentiam attulerunt hominibus, ut diceret se ante K. Martias non posse sine iniuria de provinciis Caesaris statuere, post K. Martias se non dubitaturum. Cum interrogaretur, si qui tum intercederent, dixit hoc nihil interesse utrum C. Caesar senatui dicto audiens futurus non esset an pararet qui senatum decernere non pateretur. ' quid , si,' inquit alius, 'et consul esse et exercitum habere volet?' at ille quam clementer: ' quid , si filius meus fustem mihi impingere volet?' his vocibus, ut existimarent homines Pompeio cum Caesare esse negotium, effecit. itaque iam, ut video, alteram utram ad condicionem descendere vult Caesar, ut aut maneat neque hoc anno sua ratio habeatur aut, si designari poterit, decedat.

\ciceroSection{10} Curio se contra eum totum parat. quid adsequi possit nescio ; illud video, bene sentientem, etsi nihil effecerit, cadere non posse. me tractat liberaliter Curio et mihi suo munere negotium imposuit ; nam si mihi non dedisset eas, quae ad ludos ei advectae erant Africanae, potuit supersederi ; nunc quoniam dari necesse est, velim tibi curae sit, quod a te semper petii, ut aliquid istinc bestiarum habeamus ; Sittianamque syngrapham tibi commendo. libertum Philonem istoc misi et Diogenem Graecum, quibus mandata et litteras ad te dedi. Eos tibi et rem, de qua misi, velim curae habeas ; nam quam vehementer ad me pertineat, in iis quas tibi illi reddent litteris perscripsi.
